\documentclass[10pt,a4paper]{amsart}
\usepackage{amsmath,amssymb,amsthm,graphicx}
% Retained preamble for the Stacks Project chapter snapshot in this directory.
% The chapter's own preamble.tex is a copy of the shared Stacks Project preamble; the block
% below is sliced MECHANICALLY, line for line, from it. Nothing was added, removed or
% redefined; the wrapper class, the xr-hyper/hyperref wiring and the externaldocument list
% are omitted because the assembler supplies the document class.
%
% lines taken from preamble.tex: 3, 5, 9, 11, 12, 156, 157, 158, ... (51 lines total)
%
% package coverage restored from the chapter's own preamble (2026-10-13): xy, xy-2cell

\usepackage[all]{xy}
\xyoption{2cell}
\UseAllTwocells

\newenvironment{reference}{\comment}{\endcomment}
\newenvironment{slogan}{\comment}{\endcomment}
\newenvironment{history}{\comment}{\endcomment}
\theoremstyle{plain}
\newtheorem{theorem}[subsection]{Theorem}
\newtheorem{proposition}[subsection]{Proposition}
\newtheorem{lemma}[subsection]{Lemma}
\theoremstyle{definition}
\newtheorem{definition}[subsection]{Definition}
\newtheorem{example}[subsection]{Example}
\newtheorem{exercise}[subsection]{Exercise}
\newtheorem{situation}[subsection]{Situation}
\theoremstyle{remark}
\newtheorem{remark}[subsection]{Remark}
\newtheorem{remarks}[subsection]{Remarks}
\numberwithin{equation}{subsection}
\def\lim{\mathop{\mathrm{lim}}\nolimits}
\def\colim{\mathop{\mathrm{colim}}\nolimits}
\def\Spec{\mathop{\mathrm{Spec}}}
\def\Hom{\mathop{\mathrm{Hom}}\nolimits}
\def\Ext{\mathop{\mathrm{Ext}}\nolimits}
\def\SheafHom{\mathop{\mathcal{H}\!\mathit{om}}\nolimits}
\def\SheafExt{\mathop{\mathcal{E}\!\mathit{xt}}\nolimits}
\def\Sch{\mathit{Sch}}
\def\Mor{\mathop{\mathrm{Mor}}\nolimits}
\def\Ob{\mathop{\mathrm{Ob}}\nolimits}
\def\Sh{\mathop{\mathit{Sh}}\nolimits}
\def\NL{\mathop{N\!L}\nolimits}
\def\CH{\mathop{\mathrm{CH}}\nolimits}
\def\proetale{{pro\text{-}\acute{e}tale}}
\def\etale{{\acute{e}tale}}
\def\QCoh{\mathit{QCoh}}
\def\Ker{\mathop{\mathrm{Ker}}}
\def\Im{\mathop{\mathrm{Im}}}
\def\Coker{\mathop{\mathrm{Coker}}}
\def\Coim{\mathop{\mathrm{Coim}}}
\DeclareMathSymbol{\boxtimes}{\mathbin}{AMSa}{"02}
\def\QCohstack{\mathcal{QC}\!\mathit{oh}}
\def\Cohstack{\mathcal{C}\!\mathit{oh}}
\def\Spacesstack{\mathcal{S}\!\mathit{paces}}
\def\Quotfunctor{\mathrm{Quot}}
\def\Hilbfunctor{\mathrm{Hilb}}
\def\Curvesstack{\mathcal{C}\!\mathit{urves}}
\def\Polarizedstack{\mathcal{P}\!\mathit{olarized}}
\def\Complexesstack{\mathcal{C}\!\mathit{omplexes}}
\def\Pic{\mathop{\mathrm{Pic}}\nolimits}
\def\Picardstack{\mathcal{P}\!\mathit{ic}}
\def\Picardfunctor{\mathrm{Pic}}
\def\Deformationcategory{\mathcal{D}\!\mathit{ef}}
\renewenvironment{reference}{}{}
\renewenvironment{slogan}{}{}
\renewenvironment{history}{}{}
% empty out citations inside navigation-pointer environments (no bibliography
% backend; spans themselves are cite-free - scanned 2026-10-14, lane-s3)
\renewcommand{\cite}[2][]{}
\usepackage{cleveref}
\crefname{originallemma}{original lemma}{original lemmas}
\Crefname{originallemma}{Original lemma}{Original lemmas}
\makeatletter
\newlabel{categories-definition-directed-set}{{00D3}{source}{}{sourcepointer-1}{}}
\newlabel{categories-definition-directed-set@cref}{{[originallemma][0][]00D3}{[1][0][]source}{}{}{}}
\makeatother
\crefname{originallemma}{original lemma}{original lemmas}
\Crefname{originallemma}{Original lemma}{Original lemmas}
\makeatletter
\newlabel{categories-lemma-nonempty-limit}{{086J}{source}{}{sourcepointer-2}{}}
\newlabel{categories-lemma-nonempty-limit@cref}{{[originallemma][0][]086J}{[1][0][]source}{}{}{}}
\makeatother
\crefname{originallemma}{original lemma}{original lemmas}
\Crefname{originallemma}{Original lemma}{Original lemmas}
\makeatletter
\newlabel{example-from-C-F-to-G-sets}{{0BMW}{source}{}{sourcepointer-3}{}}
\newlabel{example-from-C-F-to-G-sets@cref}{{[originallemma][0][]0BMW}{[1][0][]source}{}{}{}}
\makeatother
\crefname{originallemma}{original lemma}{original lemmas}
\Crefname{originallemma}{Original lemma}{Original lemmas}
\makeatletter
\newlabel{item-one-element}{{0BN1}{source}{}{sourcepointer-4}{}}
\newlabel{item-one-element@cref}{{[originallemma][0][]0BN1}{[1][0][]source}{}{}{}}
\makeatother
\crefname{originallemma}{original lemma}{original lemmas}
\Crefname{originallemma}{Original lemma}{Original lemmas}
\makeatletter
\newlabel{lemma-epi-mono}{{0BN0}{source}{}{sourcepointer-5}{}}
\newlabel{lemma-epi-mono@cref}{{[originallemma][0][]0BN0}{[1][0][]source}{}{}{}}
\makeatother
\crefname{originallemma}{original lemma}{original lemmas}
\Crefname{originallemma}{Original lemma}{Original lemmas}
\makeatletter
\newlabel{lemma-galois}{{0BN2}{source}{}{sourcepointer-6}{}}
\newlabel{lemma-galois@cref}{{[originallemma][0][]0BN2}{[1][0][]source}{}{}{}}
\makeatother
\title{Transitivity of the Galois group action: for a Galois category (C, F) with G = Aut(F), the action of G on F(X) is transitive for every connected X of C.}
\author{The Stacks Project Authors}
\date{stacks-project:pione.tex:local-snapshot-20261012; GNU FDL 1.2 or later}
\begin{document}
\maketitle
\noindent\textit{Editorial note.} Statement span is lines 747-755 of sources/stacks-pione/pinned-original-source.tex; proof span is lines 757-826, the complete author proof (68 body lines).\par\medskip
\noindent\textit{Editorial note.} Navigation stubs: 6 cited-but-not-included labels are supplied as non-load-bearing pointers with their pinned tags where the snapshot tag index has them.\par\medskip
\noindent\textit{Editorial note.} extra\_packages is empty on purpose: the retained preamble already loads xy with options and a second declaration clashes in the assembler wrapper.\par\medskip
\noindent\textit{Editorial note.} Text is retained verbatim, including the author's own typos and placeholders.\par\medskip
\noindent\textit{Editorial note.} Queue index 13 of queue\_pione.json (pb=3051, stmt 747-755, proof 757-826).\par\medskip
\noindent\textit{Editorial note.} pione wave-1: fresh-adjudication accept (2026-10-14); label not covered by the 85 old-session S3-pi cards.\par\medskip
\noindent\textit{Allowed standard prerequisites (editorial).}\begin{itemize}
\item Explicit prerequisite boundary. The cited inputs are prerequisites; the argument in the retained span is not replaced by any cited result.
\end{itemize}\medskip
\ifdefined\hypertarget\hypertarget{sourcepointer-1}{}\fi
\noindent\textit{Source reference.} Original lemma 00D3: sources/stacks-pione/pinned-original-source.tex, source lines 826--826. Cited by the retained statement or proof but not part of the collected spans; supplied only as a navigation pointer. The argument does not depend on its content.\par\medskip
\ifdefined\hypertarget\hypertarget{sourcepointer-2}{}\fi
\noindent\textit{Source reference.} Original lemma 086J: sources/stacks-pione/pinned-original-source.tex, source lines 826--826. Cited by the retained statement or proof but not part of the collected spans; supplied only as a navigation pointer. The argument does not depend on its content.\par\medskip
\ifdefined\hypertarget\hypertarget{sourcepointer-3}{}\fi
\noindent\textit{Source reference.} Original lemma 0BMW: sources/stacks-pione/pinned-original-source.tex, source lines 565--575. Cited by the retained statement or proof but not part of the collected spans; supplied only as a navigation pointer. The argument does not depend on its content.\par\medskip
\ifdefined\hypertarget\hypertarget{sourcepointer-4}{}\fi
\noindent\textit{Source reference.} Original lemma 0BN1: sources/stacks-pione/pinned-original-source.tex, source lines 639--647. Cited by the retained statement or proof but not part of the collected spans; supplied only as a navigation pointer. The argument does not depend on its content.\par\medskip
\ifdefined\hypertarget\hypertarget{sourcepointer-5}{}\fi
\noindent\textit{Source reference.} Original lemma 0BN0: sources/stacks-pione/pinned-original-source.tex, source lines 624--647. Cited by the retained statement or proof but not part of the collected spans; supplied only as a navigation pointer. The argument does not depend on its content.\par\medskip
\ifdefined\hypertarget\hypertarget{sourcepointer-6}{}\fi
\noindent\textit{Source reference.} Original lemma 0BN2: sources/stacks-pione/pinned-original-source.tex, source lines 705--708. Cited by the retained statement or proof but not part of the collected spans; supplied only as a navigation pointer. The argument does not depend on its content.\par\medskip
\begin{lemma}
\label{lemma-tame}
\begin{reference}
Compare with \cite[Definition 7.2.4]{BS}.
\end{reference}
Let $(\mathcal{C}, F)$ be a Galois category. Let $G = \text{Aut}(F)$
be as in Example \ref{example-from-C-F-to-G-sets}. For any connected
$X$ in $\mathcal{C}$ the action of $G$ on $F(X)$ is transitive.
\end{lemma}

\begin{proof}
We will use the results of Lemma \ref{lemma-epi-mono} without further mention.
Let $I$ be the set of isomorphism classes of Galois objects in $\mathcal{C}$.
For each $i \in I$ let $X_i$ be a representative of the isomorphism class.
Choose $\gamma_i \in F(X_i)$ for each $i \in I$.
We define a partial ordering on $I$ by setting $i \geq i'$ if
and only if there is a morphism $f_{ii'} : X_i \to X_{i'}$.
Given such a morphism we can post-compose by an automorphism
$X_{i'} \to X_{i'}$ to assure that $F(f_{ii'})(\gamma_i) = \gamma_{i'}$.
With this normalization the morphism $f_{ii'}$ is unique.
Observe that $I$ is a directed partially ordered set:
(Categories, Definition \ref{categories-definition-directed-set})
if $i_1, i_2 \in I$ there exists a Galois object $Y$ and a morphism
$Y \to X_{i_1} \times X_{i_2}$ by Lemma \ref{lemma-galois} applied
to a connected component of $X_{i_1} \times X_{i_2}$.
Then $Y \cong X_i$ for some $i \in I$ and $i \geq i_1$, $i \geq I_2$.

\medskip\noindent
We claim that the functor $F$ is isomorphic to the functor $F'$
which sends $X$ to
$$
F'(X) = \colim_I \Mor_\mathcal{C}(X_i, X)
$$
via the transformation of functors $t : F' \to F$ defined as follows:
given $f : X_i \to X$ we set $t_X(f) = F(f)(\gamma_i)$.
Using (\ref{item-one-element}) we find that $t_X$ is injective.
To show surjectivity, let $\gamma \in F(X)$. Then we can immediately
reduce to the case where $X$ is connected by the definition of
a Galois category. Then we may assume $X$ is Galois by
Lemma \ref{lemma-galois}. In this case $X$ is isomorphic to $X_i$
for some $i$ and we can choose the isomorphism $X_i \to X$ such
that $\gamma_i$ maps to $\gamma$ (by definition of Galois objects).
We conclude that $t$ is an isomorphism.

\medskip\noindent
Set $A_i = \text{Aut}(X_i)$.
We claim that for $i \geq i'$ there is a canonical map
$h_{ii'} : A_i \to A_{i'}$ such that for all $a \in A_i$
the diagram
$$
\xymatrix{
X_i \ar[d]_a \ar[r]_{f_{ii'}} & X_{i'} \ar[d]^{h_{ii'}(a)} \\
X_i \ar[r]^{f_{ii'}} & X_{i'}
}
$$
commutes. Namely, just let $h_{ii'}(a) = a' : X_{i'} \to X_{i'}$
be the unique automorphism such that
$F(a')(\gamma_{i'}) = F(f_{ii'} \circ a)(\gamma_i)$.
As before this makes the diagram commute and moreover the choice
is unique.
It follows that
$h_{i'i''} \circ h_{ii'} = h_{ii''}$
if $i \geq i' \geq i''$.
Since $F(X_i) \to F(X_{i'})$ is surjective we see that
$A_i \to A_{i'}$ is surjective.
Taking the inverse limit we obtain a group
$$
A = \lim_I A_i
$$
This is a profinite group since the automorphism groups are finite.
The map $A \to A_i$ is surjective for all $i$ by
Categories, Lemma \ref{categories-lemma-nonempty-limit}.

\medskip\noindent
Since elements of $A$ act on the inverse system $X_i$ we get an action of
$A$ (on the right) on $F'$ by pre-composing. In other words, we get
a homomorphism $A^{opp} \to G$. Since $A \to A_i$ is surjective we conclude
that $G$ acts transitively on $F(X_i)$ for all $i$. Since every connected
object is dominated by one of the $X_i$ we conclude the lemma is true.
\end{proof}

\end{document}
